\section{What the upper bound says when dimension grows}
\label{sec:dimension-oracle}
The preceding refinements make a dimension comparison possible without
hiding its entire cost in a constant. The result below fixes comparison
width and bounds the Euclidean norm of each comparison row. It concerns
the Gaussian prior; the earlier non-Gaussian extension was stated with
dimension fixed.

\begin{corollary}[A dimension-explicit shallow oracle]
\label{cor:dimension-shallow-oracle}
Fix comparison width $C$, coefficient bound $B$, accuracy $\epsilon$,
and confidence $\delta$. Bound $|c_j|\le B$ and $\|a_j\|_2\le B$
in the shallow comparison class. There is a constant $A$, depending
only on these fixed resources, such that for every $1\le d\le H$ the
Gaussian learner admits a deterministic likelihood calibration with
\[
 n\ge\exp\{A d[\sqrt H+(1-\nu)H]\}
\]
sufficient for the leading-one comparison inequality
\[
 \int R(z_w)\,d\Pi\le\inf_gR(g)+\epsilon,
 \qquad R(\bar z)\le\inf_gR(g)+\epsilon.
\]
For each $n$, one event of probability at least $1-\delta$ covers all
$N\ge C$ and all qualifying $0<\nu\le1$. Inputs may have any law on
the unit sphere, and labels lie in $[-1,1]$.
\end{corollary}
\begin{proof}
Apply Lemma~\ref{lem:hermite-encoder} at a fixed sufficiently small
accuracy. Its density norm obeys $Q\le A_0(d+1)^{a_0}$ uniformly
over the comparison class. Set $D=A_1\log(d+1)$. The dimension-uniform
schedule gives, for $H\ge A_2\log(d+1)$ and
$N\ge A_3(d+H+1)^{a_3}$, a bounded-response comparison of the desired
risk accuracy and normalized action
\[
 A_4\{1+(1-\nu)^2(H+d)+\log^2(d+1)/H\}.
\]
The original comparison output is automatically bounded by $CB$, so
all response and scalar-concentration constants are fixed here. Put
$T=1+(1-\nu)^2H+\log^2(d+1)/H$ and use
$\lambda=\ell2^{-\lceil\log_2T\rceil}$ for sufficiently small fixed
$\ell>0$. Since $d\le H$, the comparison penalty is a prescribed
small constant. The evidence cutoff satisfies $K\le A_5T$, hence
\[
 \sqrt{HK}\le A_6[\sqrt H+(1-\nu)H+\log(d+1)],
 \qquad K\le A_7(H+1).
\]
Proposition~\ref{prop:dimension-evidence} now gives the asserted
sample bound: the extra $d\log(d+H+1)$ terms are absorbed by
$A d\sqrt H$. A summable confidence allocation over the dyadic
values of $T$ gives one event for all eligible prefactors.

For $C\le N<A_3(d+H+1)^{a_3}$, the same-width weak tanh relay and
Gaussian parameter-box comparison have polynomially bounded
coefficients and inverse accuracy in $d,H$. Their normalized
comparison entropy is polynomial in $d,H$, uniformly in $\nu$.
Use a sufficiently small inverse-polynomial likelihood scale. The
finite-parameter dictionary at the resulting polynomial energy
cutoff has logarithmic cardinality polynomial in $d,H$, with degrees
depending only on the fixed comparison resources. Thus this entire
width branch needs only a polynomial sample count, which is absorbed
by the displayed exponential. The scalar concentration is applied to
one fixed risk minimizer in the compact comparison class, not to a
sample-dependent choice.

Finally, under $d\le H$, the exceptional regime
$H<A_2\log(d+1)$ contains only finitely many integer pairs $(d,H)$,
with bounds depending on $A_2$. The all-depth theorem already proved
handles these pairs at all widths. Increasing $A$ and taking their
finite maximum completes the statement.
\end{proof}

This corollary does not remove the dimension factor from the depth
exponent. For example, at canonical scaling with $d$ of order
$\sqrt H$, it gives $\exp(O(H))$ samples, whereas a lower bound of
$\exp(\Omega(\sqrt H))$ would still leave an exponential gap. More
generally, if $d$ is of order $\sqrt H+(1-\nu)H$, the sufficient
exponent is quadratic in that quantity. Thus a fixed-dimensional
$\exp(O(\sqrt H))$ upper and a varying-dimensional
$\exp(\Omega(\sqrt H))$ lower do not, by themselves, give a matched
depth dependence on continuous domains.
